\section{Anatomie des trades marquants}

Les deux figures qui suivent isolent les quatre meilleures et les quatre pires
positions. Le trait gris est le cours de clôture quotidien d'USD/JPY~; le trait
coloré relie les deux prix réellement obtenus à l'exécution.

Une convention mérite d'être posée avant de lire les chiffres. Sur l'or, le
dollar est la devise de \emph{cotation}~: le rendement en dollars d'une
position est simplement le rapport des deux cours. Sur USD/JPY, le dollar est
la devise de \emph{base}~: un notionnel de un million de dollars reste un
million de dollars quel que soit le cours, et le gain en dollars vaut
$1 - \text{prix d'entrée} / \text{prix de sortie}$. C'est cette définition
--- la seule qui reconstitue exactement le résultat du courtier --- qui est
publiée partout dans ce document. Elle est légèrement plus basse que le rapport
brut des deux cours~: le premier trade, entré à 110,761 et sorti à 139,677,
affiche $+$20,7\,\% et non $+$26,1\,\%.

\subsection{Les quatre qui ont fait le résultat}

\begin{figure}[H]
\centering
\includegraphics[width=\linewidth]{usdjpy_trades_best4}
\caption{Les quatre plus gros gains. À eux seuls, ils rapportent près du double
du résultat net du candidat.}
\end{figure}

Le trade du \textbf{24~septembre~2021}, tenu \metric{416~jours}, rapporte
$+$188\,374\,\$ --- davantage que le résultat total. Il capte le mouvement du
yen de 110,76 à 139,68, soit $+$20,7\,\% de gain de prix, auxquels s'ajoutent
\metric{25\,863\,\$} de portage encaissé sur quatorze mois. C'est le cas d'école
du momentum~: une entrée dix jours après le premier signal du run, une
tenue qui traverse toute la tendance, une sortie une fois le mouvement épuisé.

Les trois autres suivent le même schéma à plus petite échelle --- 132, 113 et
57~jours --- mais avec une différence de nature~: leur gain de prix ne dépasse
pas 3,2\,\%, et le portage y pèse entre 19 et 36\,\% du résultat. Sur ces
trades-là, le moteur n'a pas gagné en attrapant un mouvement~: il a gagné
parce qu'il a été payé pour rester en position pendant qu'il ne se passait pas
grand-chose.

\begin{insightbox}
Le quatrième panneau porte une mention particulière~: la position du
3~mars~2026 n'a pas été fermée par le signal, mais \textbf{liquidée par le
testeur} à la fin de la fenêtre, le 29~avril~2026. Son résultat,
$+$37\,343\,\$, est réel mais n'a pas été décidé par la stratégie~; il ne doit
pas être lu comme une sortie réussie.
\end{insightbox}

\subsection{Les quatre plus coûteux}

\begin{figure}[H]
\centering
\includegraphics[width=\linewidth]{usdjpy_trades_worst4}
\caption{Les quatre plus grosses pertes. Aucune n'a été coupée par le stop de
sécurité~: toutes sont des sorties de signal, sur des mouvements de change
ordinaires.}
\end{figure}

La perte la plus lourde, $-$58\,262\,\$, est celle du \textbf{3~mars~2023}~:
entrée à 135,839, sortie dix-huit jours plus tard à 132,445, soit
\metric{$-$2,6\,\%} de prix. Deux points et demi de pourcentage sur un cours de
change, c'est un mouvement banal --- et cela suffit à effacer 20\,\% du
sous-compte.

C'est le fait central de ce chapitre, et il ne se voit qu'au niveau du trade~:
\textbf{le risque de ce candidat n'est pas dans l'amplitude de ses mouvements,
il est entièrement dans son levier}. Les quatre pires pertes correspondent à
des variations de $-$1,7 à $-$2,6\,\% du change. Multipliées par un levier de
5,6 à 6,6$\times$ appliqué à un compte de l'ordre de trois cent mille dollars,
elles
deviennent des pertes à cinq chiffres.

Le corollaire est que \textbf{le stop de sécurité à 4\,\% n'a jamais servi}.
Sur l'or, il s'était déclenché une fois en quatre ans et demi, sur un décrochage
intra-séance. Sur USD/JPY, aucune séance de la fenêtre ne fait bouger le change
de 4\,\% dans le sens contraire à la position~: le garde-fou existe, il est
simplement calibré pour un actif beaucoup plus nerveux. Les trente-cinq sorties
de ce run sont toutes des sorties de signal, à l'exception de la liquidation
finale du testeur.

\begin{warningbox}
Un stop qui ne se déclenche jamais n'est pas un stop qui protège~: c'est un
stop hors d'échelle. Sur cet instrument, la perte maximale d'une position n'est
pas bornée à 4\,\% du notionnel~; elle est bornée par la fréquence de décision
du moteur, c'est-à-dire par la clôture de séance suivante. Toute intégration de
ce candidat devrait rouvrir la question du niveau du stop, sans quoi le
paramètre reste décoratif.
\end{warningbox}
